\section{The \texttt{Spot\_sample} McStas Component}
Spot sample.

\subsection*{Identification}
\begin{itemize}
  \item \textbf{Site:} 
  \item \textbf{Author:} Garrett Granroth
  \item \textbf{Origin:} Oak Ridge National Laboratory
  \item \textbf{Date:} 14.11.14
\end{itemize}

\subsection*{Description}
\begin{lstlisting}
A sample that is a series of delta functions in omega and Q.  The Q is
determined by a two theta value and an equal number of spots around the
beam center.  This component is a variation of the V sample written
by Kim Lefmann and Kristian Nielsen.  The following text comes from their
original component.
A Double-cylinder shaped incoherent scatterer (a V-sample)
No multiple scattering. Absorbtion included.
The shape of the sample may be a box with dimensions xwidth, yheight, zthick.
The area to scatter to is a disk of radius 'focus_r' situated at the target.
This target area may also be rectangular if specified focus_xw and focus_yh
or focus_aw and focus_ah, respectively in meters and degrees.
The target itself is either situated according to given coordinates (x,y,z), or
setting the relative target_index of the component to focus at (next is +1).
This target position will be set to its AT position. When targeting to centered
components, such as spheres or cylinders, define an Arm component where to
focus at.

Example: spot_sample(radius_o=0.01, h=0.05, pack = 1,
xwidth=0, yheight=0, zthick=0, Eideal=100.0,w=50.0,two_theta=25.0,n_spots=4)
\end{lstlisting}

\subsection*{Input parameters}
Parameters in \textbf{boldface} are required; the others are optional.

\begin{longtable}{p{0.22\textwidth}p{0.12\textwidth}p{0.46\textwidth}p{0.14\textwidth}}
\toprule
\textbf{Name} & \textbf{Unit} & \textbf{Description} & \textbf{Default} \\
\midrule
\endhead
radius\_o & m & Outer radius of sample in (x,z) plane & 0.01 \\
h & m & Height of sample y direction & 0.05 \\
pack & 1 & Packing factor & 1 \\
xwidth & m & horiz. dimension of sample, as a width & 0 \\
yheight & m & vert.. dimension of sample, as a height & 0 \\
zthick & m & thickness of sample & 0 \\
Eideal & meV & The presumed incident energy & 100.0 \\
w & meV & The energy transfer of the delta function & 50.0 \\
two\_theta & degrees & the scattering angle of the spot & 25.0 \\
n\_spots & 1 & The number of spots to generate symmetrically around the beam & 4 \\
\bottomrule
\end{longtable}

\subsection*{Links}
\begin{itemize}
  \item \href{run:/home/willend/willend-McCode/mcstas-comps/contrib/Spot_sample.comp}{Source code} for \texttt{Spot\_sample.comp}.
\end{itemize}
\IfFileExists{Spot_sample_static.tex}{\input{Spot_sample_static.tex}}{}